\chapter{Conclusioni e sviluppi futuri}
\label{sec:conclusion}

\section{Sintesi dei risultati}
\label{sec:concl-summary}

Questa tesi si è chiesta se un modello linguistico compatto, specializzato in un'unica decisione e impiegato come
System One, possa verificare frase per frase le risposte generate con recupero di informazioni su un portatile, e da
dove provenga l'eventuale vantaggio. GroundOne, un modello Qwen2.5 addestrato che legge i verdetti di un'intera
risposta in un solo passaggio in avanti, fornisce una risposta positiva, ma con alcune riserve.

La specializzazione funziona dove era stata indirizzata. Sull'insieme di test di RAGTruth GroundOne raggiunge
un'accuratezza bilanciata a tre classi di 0.769 con 0.5B parametri e di 0.832 con 3B, ben al di sopra degli stessi
modelli non addestrati, al di sopra degli encoder specializzati nell'ordinamento delle frasi supportate e nettamente al
di sopra di un System One generalista otto volte più grande. Su LLM-AggreFact il 3B è alla pari con il miglior encoder
di piccole dimensioni, mentre su VitaminC il vantaggio scompare: uno specialista non è un verificatore di fatti
generale. Un modello decisionale generalista della stessa dimensione, addestrato a decidere su un'ampia miscela di
compiti, conferma questo compromesso: GroundOne lo supera nettamente su RAGTruth, è alla pari su LLM-AggreFact e resta
molto indietro su VitaminC.

Il meccanismo del System One porta vantaggi su due fronti. Leggere tutti i verdetti in un solo passaggio rende la
verifica, sulle risposte lunghe, circa due volte più rapida di un giudice generativo sintetico e sei volte più rapida di
uno prolisso della stessa dimensione, con un output sempre valido; inoltre, restringere la decisione alle risposte
ammissibili è di per sé più accurato che lasciare allo stesso modello la generazione del verdetto. A parità di output, i
giudici generativi non sono in grado di fornire probabilità di qualità paragonabile: le loro probabilità separano a
malapena le decisioni corrette da quelle errate, e quelle verbalizzate sono quasi sempre del tipo tutto-o-niente.

Le probabilità di GroundOne possono essere calibrate con una temperatura e un bias stimati sui dati di sviluppo, e una
soglia scelta sugli stessi dati le trasforma in etichette coerenti sia con le probabilità sia con la decisione valutata.
L'errore di calibrazione atteso, da solo, non distingue tuttavia le probabilità utili da quelle inutili, e la capacità
discriminativa va misurata separatamente.

L'intero sistema funziona su un portatile da 16\,GB. Una risposta lunga viene verificata in 0.41 secondi e 5.5 joule
con il modello 0.5B, e in 2.57 secondi e 33 joule con il 3B. Il risparmio di energia rispetto a un giudice è reale ma
inferiore al risparmio di tempo, perché un singolo prefill lungo assorbe più potenza della decodifica.

I limiti principali risiedono nella ristrettezza dei dati di addestramento, a causa della quale le contraddizioni al di
fuori del data-to-text e le allucinazioni nascoste in frasi copiate vengono apprese solo in parte, e in una procedura di
addestramento non del tutto robusta rispetto al seed casuale.

\section{Sviluppi futuri}
\label{sec:concl-future}

Da questi limiti discendono direttamente diverse direzioni di lavoro. La procedura di addestramento andrebbe resa
robusta abilitando la correzione del bias e il clipping del gradiente, registrando le norme dei gradienti e ripetendo le
esecuzioni con più seed, anche per il modello 3B. Dati di addestramento con contraddizioni più varie, con supporto
parafrasato e con modifiche minime del tipo contenuto in VitaminC dovrebbero affrontare sia il richiamo debole al di
fuori del data-to-text sia il punto cieco sulle frasi copiate. Addestrare la lettura dei logit (readout) sulla lettera che il modello
scriverebbe spontaneamente dopo il marcatore potrebbe allineare readout e generazione, con effetti sull'accuratezza che
restano da misurare.

Sul versante dei sistemi, le misure ancora mancanti per il pilastro della località sono la memoria occupata quando il
verificatore opera accanto a un generatore e a un sistema di recupero, e una valutazione completa con la rete
disattivata. Un confronto equo della latenza con gli encoder richiede di eseguirli in batch, mentre quello con il
System One generalista richiede di confrontare entrambi i sistemi attraverso la stessa interfaccia HTTP. La verifica
potrebbe inoltre avvenire durante la generazione della risposta, segnalando ogni frase non appena è completa, cosa che
la soluzione a marcatori rende naturale. Infine, un assistente locale costruito sull'interfaccia del servizio
consentirebbe agli utenti di scegliere la regola e la soglia più adatte ai propri documenti, ed è proprio in questo
contesto che le probabilità calibrate mostrano il loro valore.
